\section{The conditions in full and the component statements}
\label{sec:moved-main-2026-10-13}

The main text states its most important results in full and the others in words.
This section holds the precise statements of the others, with their discussions, and
the precise forms of the conditions that the main text names in words. Every
cross-reference to them resolves here.

\subsection{The bounds of the class and the conditions in full}
\label{sec:moved-main-2026-10-13:conditions}

\central{Except where a statement names the quadratic--linear family of
\Cref{sec:model:family}, the results use no functional form for the transition
or for either payoff.} They hold on the \emph{general class} of economies
defined by \textup{(F1)--(F5)}, by bounds on the transition's slopes and
curvature, and by bounds on the payoffs' rates. Subscripts $\mathrm{lo}$ and $\mathrm{hi}$ mark a
lower and an upper bound. The transition is constrained
only through its chords, with no derivative assumed. Its \emph{slope bounds}
bound the slope, or difference quotient, of every chord, with the state slope
in $[g_{\mathrm{lo}},g_{\mathrm{hi}}]$ at every sanction and the sanction slope
in $[d_{\mathrm{lo}},d_{\mathrm{hi}}]$, $d_{\mathrm{lo}}>0$. The transition
also satisfies a joint semiconcavity bound. Some results read two further
curvature bounds (\Cref{app:proofs:placable}). One is a semiconvexity modulus
$M_{\mathrm{cx}}$, which limits how far the transition can bend in the state.
The other is a two-sided bound $M_{\mathrm{mix}}$ on mixed differences, which
limits how much the state changes a sanction's effect on the next state. These
two moduli are read only by results that bound the curvature of the sender's
value from above and below, and by one route from bounds on the target's policy.
So is the sender's semiconvexity modulus $\tau_x$ below. Among those results is
the clause of \Cref{prop:placable-family} that uses the increasing-differences
condition \eqref{eq:rally-margin}.
Membership of the class does not require these three moduli.

\central{The rate bounds fall into four groups.} \emph{The sender's stake}: its
payoff rises with the state at a rate of at least $w_{\mathrm{lo}}>0$, and at
full non-compliance at most $w_{\mathrm{hi}}$, with a semiconvexity modulus
$\tau_x\ge0$ in the state there. \emph{The sender's sanction choice}: the
curvature of its sanction cost lies between $q$ and $\bar Q$, and $q>0$ follows
from the concavity condition of \Cref{sec:model:conditions}. At full non-compliance
its marginal return at the first unit of sanction lies in
$[u_{\mathrm{lo}},u_{\mathrm{hi}}]$, with $u_{\mathrm{lo}}>0$. A higher state
raises that return at a rate of at least $r_{\mathrm{lo}}\ge0$, the sender's
\emph{state--sanction complementarity}. Compliance lowers it at a substitution rate
of at least $k_{\mathrm{lo}}>0$, and some results also assume an \emph{upper
substitution rate} $k_{\mathrm{hi}}$. \emph{The target's costs}: compliance costs
between $\ell_{\mathrm{lo}}>0$ and $\ell_{\mathrm{hi}}$ per unit, and the
sanction's marginal pain lies in $[p_{\mathrm{lo}},c]$, with $p_{\mathrm{lo}}>0$.
\emph{The target's stake}: its payoff rises with the state at a rate between
$v_{\mathrm{lo}}+n_{\mathrm{lo}}s$ and $v_{\mathrm{hi}}+n_{\mathrm{hi}}s$ at
sanction $s$, with $v_{\mathrm{lo}}>0$ and $n_{\mathrm{lo}}\ge0$.

\central{The theorems use five constants built from these payoffs and bounds.}
Two of them define the entrenchment condition and its strong form. Write
$K_S=w_{\mathrm{hi}}/(1-\beta_S g_{\mathrm{hi}})$ for the Lipschitz constant of
the sender's value in the state, and
$A=u_{\mathrm{hi}}+\beta_S K_S d_{\mathrm{hi}}$ for an upper bound on the
sender's dynamic marginal return to sanctioning at full non-compliance. Write
$\underline s=\sigM(0,0)$ for the sanction that maximizes the sender's
per-period payoff against full non-compliance at $x=0$ (the myopic best response
$\sigM$ is defined in \Cref{sec:model:myopic}). Finally, write
\begin{equation}
g_T^0=\frac{v_{\mathrm{lo}}}{1-\beta_T g_{\mathrm{lo}}},
\qquad
g_T=\frac{v_{\mathrm{lo}}+n_{\mathrm{lo}}\underline s}{1-\beta_T g_{\mathrm{lo}}}
\label{eq:target-slopes}
\end{equation}
for a lower state slope of the target's value and its refined version, which
uses $\underline s$ as a floor on the sender's response at non-compliance: the target's
own stake per period, discounted over the persistence of the state. The two
slopes give the \emph{entrenchment condition} (\Cref{ass:loss-offset})
\begin{equation}
c\le\beta_T\,d_{\mathrm{lo}}\,g_T ,
\label{eq:loss-offset}
\end{equation}
under which the discounted value of the entrenchment a unit of sanction builds
covers the incumbent's current pain from it, and the \emph{strong entrenchment
condition} $c\le\beta_Td_{\mathrm{lo}}g_T^0$ (\Cref{sec:results:security}). The
strong form uses the unrefined slope, so it is more demanding. On the sender's
side, write $g_S=w_{\mathrm{lo}}/(1-\beta_S g_{\mathrm{lo}})$ for the least rate at
which a sender value in the invariant set rises with the state. Under interiority
the forward-looking sender's sanction at non-compliance exceeds the myopic sanction
$\sigM(x,0)$ at every state by at least
\begin{equation}
\frac{\beta_S\,g_S\,d_{\mathrm{lo}}}{\bar Q}>0 ,
\label{eq:sender-wedge}
\end{equation}
the premium bound of \Cref{sec:results:wedge}: the numerator is the least that a unit
of sanction adds to the sender's discounted value through the next state, and
$\bar Q$ converts that marginal gain into a sanction increase.

Two remarks of the earlier main text belong here. The formalization maintains
\textup{(F2)} and \textup{(F3)} on all of $\mathbb{R}$ at no cost in generality: a
transition satisfying \textup{(F1)--(F4)} on the square extends globally, unchanged
on the square, with the same properties, and a machine-checked constructor builds
the formal economy from square-defined data (\Cref{app:moved-ad}). For the
closed-form constants of \Cref{app:closedforms} the payoffs are also twice
continuously differentiable on a neighborhood of $[0,1]^3$; every theorem in the
formalization holds without this (\Cref{app:moved-ad}).

\central{Of the conditions listed below, only the concavity condition is part
of membership of the class.} Every other condition is a
hypothesis of the results that read it. The concavity condition says that the
curvature $q$ of the sender's sanction cost exceeds the curvature that the
continuation term can add. This curvature is computed with a curvature bound
$\sigma_V$ on the sender's value that the Bellman operator preserves. Such a
bound exists exactly when an explicit inequality on the primitives holds
(\Cref{app:proofs:placable}). Under this condition the sender's dynamic
objective is strictly concave in the sanction. Interiority $A<q$ implies the
myopic interiority $u_{\mathrm{hi}}<q$, under which a myopic sender never
sanctions fully, and some results read only that (\Cref{thm:regime-separation}
and \Cref{prop:target-only}(iv)). The lifting condition $A\le k_{\mathrm{lo}}$ implies
$u_{\mathrm{hi}}<k_{\mathrm{lo}}$, under which full compliance ends a myopic
sender's sanction. That is all the myopic equilibrium reads of it. The myopic
relief condition $\ell_{\mathrm{hi}}\bar Q<k_{\mathrm{lo}}p_{\mathrm{lo}}$ is a
one-period margin with the sender's response substituted. The reversal toward
non-compliance is a continuation-value result. Five bounds are not required of every
economy in the class. For the target, they are the upper bound
$\ell_{\mathrm{hi}}$ on its compliance cost, the lower bound $p_{\mathrm{lo}}$
on its pain and the upper bound $v_{\mathrm{hi}}+n_{\mathrm{hi}}s$ on its stake.
For the sender, they are the bound away from non-compliance on its state--sanction complementarity and the upper
substitution rate $k_{\mathrm{hi}}$.
\Cref{prop:branch-invariance,prop:relief-participation,prop:monotone-uniqueness}
hold without them, and the other results that use $k_{\mathrm{hi}}$ take it as
a hypothesis. Every example economy in the class has an upper substitution rate and the moduli
$M_{\mathrm{cx}}$, $M_{\mathrm{mix}}$ and $\tau_x$.
\central{Each theorem-specific condition does one economic job and is read by
the results named with it.}
\begin{itemize}
\item \emph{Concavity condition} (membership): a well-defined dynamic best
response and controlled continuation values.
\item \emph{Interiority} $A<q$: an interior dynamic best response at non-compliance.
The myopic and target-only cases, and the result that the sender's
forward-looking incentive changes behavior, read only the myopic interiority
$u_{\mathrm{hi}}<q$ it implies.
\item \emph{Myopic relief condition}
$\ell_{\mathrm{hi}}\bar Q<k_{\mathrm{lo}}p_{\mathrm{lo}}$: at the margin, a unit
of compliance costs at most $\ell_{\mathrm{hi}}$ and buys at least
$k_{\mathrm{lo}}/\bar Q$ units of sanction reduction worth at least
$p_{\mathrm{lo}}$ each, so the myopic target accommodates. It is read by the
myopic case and by the threshold discount factor of
\Cref{prop:collapse-threshold}, and implied by the forward-looking relief condition
\textup{(RL)} of \Cref{prop:sender-only}.
\item \emph{Lifting condition} $A\le k_{\mathrm{lo}}$: full compliance drives the
sanction to zero. It is read by cases 00 and 01 (myopic target) among other
results, case 00 only through the inequality $u_{\mathrm{hi}}<k_{\mathrm{lo}}$ it
implies (full compliance ends a myopic sender's sanction).
\item \emph{Compliance neutrality} \eqref{eq:zero-sanction-neutrality}: it
determines the sender's value along relief before the target policy is known,
and only the results that construct relief read it.
\item \emph{Entrenchment condition} $c\le\beta_Td_{\mathrm{lo}}g_T$: the
discounted value of the entrenchment a unit of sanction builds covers the
incumbent's current pain from it. It is read by cases 10 and 11
(forward-looking target) among other results; its strong form
$c\le\beta_Td_{\mathrm{lo}}g_T^0$ is read by the minmax-payoff and uniqueness
results.
\item \emph{Increasing-differences condition} \eqref{eq:rally-margin}: a
nondecreasing sender response at non-compliance, read by the constructed full-non-compliance
equilibrium of \Cref{prop:placable-family}.
\item $\ell_{\mathrm{hi}}<\underline s\,p_{\mathrm{lo}}$: full compliance beats
full non-compliance in the one-period problem. It is read by the clause of
\Cref{prop:placable-family} that the forward-looking incentives change behavior
and by the two results of \Cref{app:proofs:placable} that derive the forward-looking
relief condition \textup{(RL)} from it.
\end{itemize}

\subsection{The components of the four-case theorem}
\label{sec:moved-main-2026-10-13:components}

\Cref{prop:placable-family} constructs the full-non-compliance equilibrium used in
case $11$; \Cref{thm:regime-separation} and
\Cref{prop:target-only,prop:sender-only} are the components of
\Cref{thm:four-cells} in the main text.

\begin{proposition}[The sender's response under interiority and the lifting condition]
\label{prop:placable-family}
For every economy in the class satisfying interiority $A<q$, the sender's value
against full non-compliance is the
Banach fixed point of its Bellman operator on an invariant set of Lipschitz,
lower-sloped, curvature-bounded value functions, and its dynamic best response
never reaches the ceiling:
every maximizer of $J_S$ lies in $[0,A/q]$ with $A/q<1$, and
at full non-compliance it is strictly interior and
strictly above the myopic sanction $\sigM(x,0)$ ($(u_0+rx)/q$ on the
family); under the lifting condition it
is zero at full compliance. If in addition the transition has a state-section
semiconvexity modulus $M_{\mathrm{cx}}$ and a bound $M_{\mathrm{mix}}$ on mixed
differences, and the sender's payoff a semiconvexity modulus $\tau_x$ in the state
(\Cref{sec:moved-main-2026-10-13:conditions}), and the economy satisfies the entrenchment condition and the
increasing-differences condition
\begin{equation}
\beta_S\bigl(\sigma_C g_{\mathrm{hi}} d_{\mathrm{hi}}
+(K_S+2\sigma_C)M_{\mathrm{mix}}\bigr)\le r_{\mathrm{lo}},
\qquad
\sigma_C=\frac{\tau_x+\beta_S K_S M_{\mathrm{cx}}}{1-\beta_S g_{\mathrm{hi}}^2},
\label{eq:rally-margin}
\end{equation}
then the response at non-compliance is non-decreasing in the state, the
target's value is the Banach fixed point of its Bellman operator against
that response, full non-compliance is the target's unique optimum at every
state, and $(\Vt,\Vs,\alpha\equiv0)\in\mathcal E$. At that equilibrium, at
every state, no maximizer of $J_S$ is a stage maximizer; and
when $\ell_{\mathrm{hi}}<\underline s\,p_{\mathrm{lo}}$, so that full compliance
beats full non-compliance in the one-period problem, no
forward-looking target maximizer is a stage maximizer either.
\end{proposition}

\central{The increasing-differences condition \eqref{eq:rally-margin} says that
the stage complementarity between state and sanction covers two forward-looking
forces that can make the sender's response fall with the state.} They are the
curvature of the sender's continuation value (the $\sigma_C$ term) and the
dependence on the state of the sanction's effect on the next state (the
$M_{\mathrm{mix}}$ term). Both transition terms, $M_{\mathrm{cx}}$ and
$M_{\mathrm{mix}}$, vanish for every additively separable transition with convex
state sections, such as the partial-adjustment law
(\Cref{sec:model:examples}). With
a sender payoff convex in the state ($\tau_x=0$, as on the family), the
$\sigma_C$ term vanishes too, and the condition reads $0\le r_{\mathrm{lo}}$.
The construction uses \eqref{eq:rally-margin} only to make the response at
non-compliance non-decreasing. The appendix states
\Cref{thm:regime-separation,thm:four-cells}, \Cref{prop:target-only}(iv), the
long-run gap of \Cref{prop:steady-gap} and the welfare results of
\Cref{sec:welfare:nonidentification} under that monotonicity instead. It is
strictly weaker than \eqref{eq:rally-margin} with the supplied moduli: it holds
at a member that meets the condition with a tighter modulus but not with the
supplied ones
(\Cref{app:proofs:placable}). We state \eqref{eq:rally-margin} because, unlike that
monotonicity, it can be checked from the primitives and the supplied moduli.

\central{Relief and non-compliance are constructed by different routes, which is why
the results use different conditions.} At non-compliance, the sender's value against
full non-compliance is constructed first and the target's value against the resulting
response next. Only then do the entrenchment and increasing-differences
conditions verify that non-compliance is the target's optimum, so that the
construction against a fixed full-non-compliance policy becomes an equilibrium. In
cases $00$ and $01$, where the myopic target accommodates, compliance
neutrality pins down the sender's value along relief, and the target's optimal
compliance is the least that makes the selected sanction zero. The minmax bound
on the target's value and the uniqueness of full non-compliance among equilibria with
a non-decreasing sanction rule therefore hold without the conditions that
construct relief. Conversely, a relief equilibrium requires solving
its own continuation problem. The myopic and forward-looking cases are \Cref{prop:myopic-benchmark} and
\Cref{thm:regime-separation} of the main text; the target-only and sender-only cases
follow (\Cref{prop:target-only,prop:sender-only}).

The separation of the myopic and the forward-looking cases is in the main text: the myopic case is \Cref{prop:myopic-benchmark},
and the forward-looking case, with the comparison of the two paths, is
\Cref{thm:regime-separation}, both under the one-period forms of interiority and
the lifting condition.

\central{The two main-text results are statements about the constructed
equilibria, and their parts rest on different conditions.} The forward-looking
equilibrium of \Cref{thm:regime-separation} is constructed under the increasing-differences condition. Given it, the lower bound on the build-up of the state
and the forward-looking path lying above the myopic one use only the transition's
monotonicity and sanction-slope bounds. The convergence also uses the response
monotonicity that the condition delivers. The convergence to zero in
\Cref{prop:myopic-benchmark} uses the sanction-slope bounds alone. The limit in
\Cref{thm:regime-separation} may depend on the initial state unless the
equilibrium law of motion is a contraction (\Cref{prop:activation-not-hysteresis}).
Its bounds $\underline s<x^{11}_\infty\le A/q$ have closed forms on the family
(\Cref{sec:model:examples}), and $A/q<1$ under interiority. The comparison changes
both discount factors together, holding the other primitives fixed.

\begin{proposition}[Target-only case]
\label{prop:target-only}
For every economy in the class satisfying the entrenchment condition, with
$\beta_S=0$ and $\beta_T$ as given:
\begin{enumerate}[label=(\roman*)]
\item the sender's sanction rule is the myopic sanction $\sigM(x,0)$
($(u_0+rx)/q$ on the family under interiority) and its value is the stage
value there;
\item full non-compliance is the target's unique optimum at every state, and its
value is the fixed point of its Bellman operator against the myopic best response;
the pair with $\alpha\equiv0$ is an equilibrium in $\mathcal E$ of this
case;
\item from every initial state $x_t\ge\underline s\bigl(1-(1-d_{\mathrm{lo}})^t\bigr)$
and $x_t$ converges to a fixed point $x^{10}_\infty\ge\underline s$ of the myopic sanction rule;
on the family under interiority the fixed point is unique, $u_0/(q-r)$, and
$\bigl|x_t-u_0/(q-r)\bigr|\le\bigl(1-d_{\mathrm{lo}}(1-r/q)\bigr)^t\,\bigl|x_0-u_0/(q-r)\bigr|$
($7/12$ per period at the partial-adjustment example);
\item under the increasing-differences condition \eqref{eq:rally-margin} and
$u_{\mathrm{hi}}<q$ (a myopic sender never sanctions fully), the limit $x^{10}_\infty$ from a
start is strictly below the limit $x^{11}_\infty$ of the forward-looking equilibrium of
\Cref{thm:regime-separation} from the same start.
\end{enumerate}
\end{proposition}

\central{The target's forward-looking incentive alone therefore switches it from
partial compliance to non-compliance, and the sender's forward-looking incentive then
raises the steady state.} The entrenchment condition does that work against the
myopic best response as against the dynamic one. The steady state rises
because the forward-looking sanction exceeds the myopic sanction at every
state. On the family, the rate in (iii) bounds each path's distance to the
steady state at each date. Across economies of the family that share $u_0$ and
$q$, a larger sender state--sanction cross term $r$ strictly raises the myopic
sanction at every positive state where it is below full sanction. It strictly raises the target-only steady
state too when the economy with the larger $r$ satisfies interiority
(\Cref{app:moved-ad}).

\begin{proposition}[Sender-only case]
\label{prop:sender-only}
For every economy in the class satisfying interiority, the lifting condition and
compliance neutrality whose transition also has
semiconvexity in the sanction, condition \textup{(T6)} of
\Cref{app:proofs:placable}, satisfying
the forward-looking relief condition \textup{(RL)} there, with $\beta_T=0$ and $\beta_S$ as
given:
\begin{enumerate}[label=(\roman*)]
\item the sender's value is the value of the zero-sanction path,
$\sum_{t\ge0}\beta_S^t\,\Us(y_t,0,0)$ with $y_0=x$ and $y_{t+1}=\F(y_t,0)$
($\Us(x,0,0)=wx$ on the
family); it is constructed before
the target's policy is determined, because along the constructed equilibrium
the sanction is zero and at sanction zero the sender's stage payoff
(by compliance neutrality) and the transition do not depend on compliance;
\item the target's unique optimum at every state is
$a^{01}(x)$, the least compliance at which the sender's dynamic best
response vanishes; it is continuous in the state, lies strictly above
$a^{00}(x)$, the compliance that ends a myopic sender's sanction, and with it the pair
$(\Ut(x,a^{01}(x),0),\,V_S,\,a^{01})$ is an equilibrium in $\mathcal E$ of this
case ($\Ut(x,a^{01}(x),0)=v\,x-\ell a^{01}(x)$ on the family);
\item the equilibrium sanction is zero and from every initial state
$x_t\le(1-d_{\mathrm{lo}})^tx_0\to0$.
\end{enumerate}
\end{proposition}

\central{The forward-looking sender is more demanding than the myopic one, so
the myopic target pays more for relief, but it still buys it.} Because of the
foresight premium, the forward-looking sender still sanctions at the compliance
$a^{00}(x)$. The forward-looking relief condition (RL),
$\ell_{\mathrm{hi}}M_{\mathrm{tot}}<k_{\mathrm{lo}}p_{\mathrm{lo}}$, says that a unit of compliance
avoids more sanction pain, valued at the lowest price $p_{\mathrm{lo}}$, than it
costs. Here $M_{\mathrm{tot}}$ bounds the curvature of the sender's dynamic objective
in the sanction (\Cref{app:proofs:placable,app:moved-ah}). Relief can require
full compliance (\Cref{app:moved-ad}).

\subsection{The strict comparison at a fixed target policy}
\label{sec:moved-main-2026-10-13:strict}

\begin{proposition}[Strict comparison at a fixed target policy]
\label{prop:strict-dominance}
Fix $x$ and a target action $a$, write $\Gamma(s)=\Vs(\F(x,s))$ for the
sender's continuation value at the successor state $\F(x,s)$, and let $m$ be the greatest stage maximizer of
$U_S(x,a,\cdot)$ on $[0,1]$, which exists by continuity. If no feasible
sanction below $m$ carries a larger continuation value, $\Gamma(s)\le\Gamma(m)$ for every $s\in[0,m]$, and
there is one sanction $s'$ at which the stage loss measured from $m$ is
outweighed by the discounted continuation gain,
\[
U_S(x,a,m)-U_S(x,a,s')<\beta_S\bigl(\Gamma(s')-\Gamma(m)\bigr),
\]
then every stage maximizer lies weakly below $m$ and strictly below every
maximizer of $J_S$.
\end{proposition}

\central{\Cref{prop:strict-dominance} needs no concavity, differentiability,
uniqueness or positive discount, and $\Gamma$ need not be monotone: a
non-decreasing $\Gamma$ is one way to meet the comparison.} The proposition
orders two sets of maximizers, and the stage loss must be measured from the
greatest stage maximizer. At an economy in
\Cref{app:moved-ah}, a loss measured from a lower tied stage maximizer orders
only that one. The proof closes \Cref{app:proof-dominance}.

\central{Compared with the myopic equilibrium at their own target policies, the
dynamic equilibrium sanctions at least as much at every state when its sender value is non-decreasing in the state, its target
complies no more and compliance placates the sender.} It does so through two
separate effects.
Under the compliance ordering $\alpha(x)\le\alphaM(x)$,
\[
\sigM(x,\alphaM(x))
\le \sigM(x,\alpha(x))
\le \sigDynTop_{\Vs}(x,\alpha(x)).
\]
The first inequality is the indirect effect of the target's lower compliance on
the myopic best response. The second is the direct effect of the sender's foresight
at that compliance (\Cref{thm:dominance}). The compliance ordering is an
explicit condition, met by the example economies. Under the general transition, whole sanction paths are ordered under other conditions, without the compliance ordering
or any condition on how the target ranks sanction levels. This is so
for an equilibrium constructed from stated invariant
value sets, when the myopic sanction stays below a ceiling on the high political
basin and the myopic rule preserves that basin
(\Cref{cor:pathwise-dominance}).

\subsection{The target's minmax payoff and uniqueness among monotone rules}
\label{sec:moved-main-2026-10-13:security}

\central{The benchmark for the target's security is permanent relief: zero
sanctions at every date, with the target maintaining non-compliance throughout.} Let $y_0=x_0$ and
$y_{t+1}=\F(y_t,0)$ be the state path under permanent zero sanctions, and define
the permanent-relief payoff
\begin{equation}
\underline{V}_T(x_0)=\sum_{t=0}^{\infty}\beta_T^t\,\Ut(y_t,0,0).
\label{eq:relief-benchmark}
\end{equation}

\begin{proposition}[The target's minmax payoff]
\label{prop:security}
Take any economy, whatever the sender's objective, in which the transition's
zero-sanction section rises in the state at rate at least $g_{\mathrm{lo}}$, a
sanction $s$ raises the successor state by at least $d_{\mathrm{lo}}s$ against
zero sanctions, and the target gains at least $v_{\mathrm{lo}}\ge0$ per unit of
state at full non-compliance and zero sanctions, loses at most $cs$ to a sanction $s$
at full non-compliance, and gains nothing from compliance against a zero sanction
(used only in (i) and (ii)). Under the strong entrenchment condition
$c\le\beta_Td_{\mathrm{lo}}g_T^0$, and from every initial state $x_0$:
\begin{enumerate}[label=(\roman*)]
\item $\underline{V}_T(x_0)$ is the lowest payoff to which a feasible pure
history-dependent sanction strategy can hold the target against every
feasible pure history-dependent compliance strategy; permanent relief attains it;
\item $\underline{V}_T(x_0)$ is the highest payoff that the target can guarantee
itself with a feasible pure history-dependent compliance strategy against every
feasible pure history-dependent sanction strategy; full non-compliance at every date
attains it;
\item every equilibrium in $\mathcal E$ satisfies $\Vt(x_0)\ge \underline{V}_T(x_0)$.
\end{enumerate}
\end{proposition}

\central{So $\underline{V}_T$ is the target's minmax payoff in pure strategies,
which by (ii) is also the most it can guarantee itself, and the minmax over
sanction sequences fixed in advance is its special case.} Every economy in the
class has the five properties the proposition reads, from the bounds
$g_{\mathrm{lo}}$, $d_{\mathrm{lo}}$, $v_{\mathrm{lo}}$ and $c$ and the positive
compliance cost. Inequality (iii) holds at every equilibrium in $\mathcal E$,
beyond the constructed equilibrium and the monotone-rule equilibria of
\Cref{prop:identification}.

\central{The bound holds because, under the strong entrenchment condition, the
entrenchment that sanctions build is worth at least as much to the target as
their pain.} Under any sanction sequence, the lower bounds on the transition's
slopes make the state's excess over the zero-sanction state a discounted
accumulation of past sanctions. Valuing it at least at its lower state-valuation
rate, against a pain of at most $c$ per unit, the target gains from full
non-compliance at every date at least
$\bigl(\beta_Td_{\mathrm{lo}}v_{\mathrm{lo}}/(1-\beta_Tg_{\mathrm{lo}})-c\bigr)\sum_t\beta_T^ts_t$
over permanent relief. The strong entrenchment condition makes this nonnegative.
Against permanent relief, compliance only incurs its cost, so full non-compliance
attains the upper bound. With deterministic play and perfect monitoring,
date-indexed actions reproduce the path of any pair of history-dependent plans,
so both bounds extend to history-dependent strategies. Equilibrium optimality gives the target at least the payoff of
deviating to full non-compliance at every date (\Cref{app:proofs:scope}).
Politically, under the strong entrenchment condition relief is the harshest
treatment the sender can impose, so threats cannot coerce the target on the
non-compliance side.

\central{The same strong entrenchment condition makes full non-compliance, with the
associated values, the only possible equilibrium with a non-decreasing sanction
rule, and excludes a
sanction rule equal to zero at every state.}

\begin{proposition}[Uniqueness of non-compliance among monotone rules]
\label{prop:identification}
On an economy in the class satisfying the strong entrenchment condition
$c\le\beta_Td_{\mathrm{lo}}g_T^0$:
\begin{enumerate}[label=(\roman*)]
\item every equilibrium in $\mathcal E$ whose sanction rule is
non-decreasing in the state has full non-compliance at every state, the constructed
sender and target values, and the constructed response as its sanction rule;
\item no equilibrium in $\mathcal E$ has a sanction rule equal to zero at
every state.
\end{enumerate}
\end{proposition}

\central{Both clauses of \Cref{prop:identification} concern non-decreasing
sanction rules: (i) pins down every equilibrium with such a rule, and (ii)
excludes the rule equal to zero at every state.} Clause (i) requires no monotonicity of the target's policy.
Its proof rules out positive compliance in the class through a lower state slope
that the equilibrium value equation gives the target's value
(\Cref{prop:monotone-uniqueness}). Clause (ii) covers every equilibrium: a
globally zero rule is non-decreasing, so it would fall within (i), which implies
a positive response at full non-compliance. Under the increasing-differences
condition the constructed equilibrium belongs to the class of (i), which is
therefore not empty.

\subsection{Relief, non-compliance and the threshold discount factor in general}
\label{sec:moved-main-2026-10-13:bracket}

\central{For the forward-looking target, three inequalities on its discount
factor each give one conclusion (the formulas are in \Cref{prop:bracket}).}
(i)~\emph{Sufficient condition for non-compliance.} The most the target can save by buying
off a unit of sanction, its pain $c$ less the least compliance cost of that
unit, is below the discounted entrenchment the unit builds. That entrenchment is valued at a lower
slope that the target's value carries when the sender's response is monotone.
The slope is $g_T$ under the entrenchment condition and, in general, under two further
conditions, a slope at most $g_T$. Then full non-compliance is the target's unique
optimum at every state, and with the increasing-differences condition the
constructed full-non-compliance equilibrium exists. (ii)~\emph{Sufficient condition
for relief.} The most a unit of compliance costs, $\ell_{\mathrm{hi}}$, is below
the least value of what it buys. This is $k_{\mathrm{lo}}/M_{\mathrm{tot}}$ units of
sanction reduction, each worth the pain $p_{\mathrm{lo}}$ less the most
discounted entrenchment it can cost in the relief equilibrium. Then, under
interiority, the lifting condition and compliance neutrality, on economies in the class
with semiconvexity in the sanction \textup{(T6)}, the relief equilibrium, in
which the target complies at $a^{01}(x)$, exists. For existence the weak forms of
(i) and (ii) suffice. The strict forms add uniqueness of the target's optimum.
(iii)~\emph{An inequality that rules out non-compliance.} It compares as in (ii),
measuring entrenchment by the Lipschitz constant of the target's value at the
full-non-compliance profile. With interiority, the lifting condition and the target's discount
factor times a bound on the Lipschitz constant of the state map at full
non-compliance below one, it makes full non-compliance not a best reply at that profile's
values. Since every slope in (i) is at most $g_T$, each sufficient condition for
non-compliance implies the inequality for non-compliance at $g_T$, which excludes (ii) and,
under the same bound on the state map, (iii). The sufficient condition for
non-compliance and the inequality at $g_T$ pass to a more patient target and (ii) to a
less patient one. In strict form, the inequality at $g_T$ and (ii) are open
conditions on $\beta_T$.

\central{In general these inequalities leave a gap between the two sufficient
conditions.} Above every
discount factor at which the sufficient condition for relief holds and below
every discount factor at which the sufficient condition for non-compliance holds,
neither applies. They determine no equilibrium there. The gap can reflect
conservative bounds as well as genuine differences in behavior. Each end has
an exact counterpart, a threshold necessary and sufficient for existence, under
the conditions of \Cref{prop:exact-thresholds}. At the relief end these are \textup{(T6)},
interiority, the lifting condition, compliance neutrality and a target state valuation
at zero sanction that covers the rise of its compliance cost along the relief
policy. At the non-compliance end they are a non-decreasing sender response at full non-compliance
and the slope condition of \Cref{prop:bracket}(i) at slope zero. The discount
factors with an equilibrium without sanctions then form an interval from zero up
to a largest one, and those with a full-non-compliance equilibrium an interval from a
smallest discount factor up to one. At the partial-adjustment example with $\beta_T$ varied
(\Cref{sec:model:examples}), where an exact evaluation finds the gap open
(\Cref{app:proofs:threshold}), both apply, and both exact thresholds lie
strictly inside the gap. There, at one discount factor between them, a continuum of
equilibria with compliance rising in the state exists (\Cref{prop:affine-window}).

\central{Nine equalities, the linear--quadratic special case of the main text, close this gap: with them the sufficient condition for
relief and the sufficient condition for non-compliance meet at a single discount
factor.} They make the transition affine on
the square by hypothesis (a constant sanction slope, with $\F(x,x)=x$); they are listed in
the proof of \Cref{prop:collapse-threshold} (\Cref{app:proofs:scope}). They also set two
curvature bounds on the sender's side to zero, rule out state--sanction cross
terms and make five marginal rates constant. The five
further equalities of \Cref{prop:bracket}(iii), tight bounds on the transition's
state slope and four zero transition moduli, concern the bounds supplied
for the economy and leave the game unchanged (\Cref{app:proofs:scope}). In symbols, the
transition's constant sanction slope is $\delta$ (so that $\F(x,s)=(1-\delta)x+\delta s$), the
constant curvature of the sender's cost of sanctioning is $q$, and the constant
substitution rate, compliance cost, sanction pain and state valuation of the target are
$k$, $\ell$, $c$ and $v$. Write $D=c-\ell q/k$ for the target's net cost of a sanction, its pain less the
compliance cost of ending it. The threshold discount
factor of \Cref{prop:collapse-threshold} is
\begin{equation}
\beta_T^{*}=\frac{D}{\delta v+(1-\delta)D},
\label{eq:target-threshold}
\end{equation}
which rises with $D$ and falls with $v$; under the nine equalities the myopic relief
condition reads $\ell q<kc$, that is, $D>0$.

\central{The non-existence direction ranges over every equilibrium in
$\mathcal E$: below $\beta_T^*$ no equilibrium has full non-compliance.} At
$\beta_T=\beta_T^*$ the equilibrium exists,
because existence passes to the limit in $\beta_T$
(\Cref{prop:defiance-closed}). Two example economies lie on the two sides
of the threshold (\Cref{app:proofs:scope}). When
$\beta_T^*\ge1$, which holds exactly when $D\ge v$, no discount factor
supports full non-compliance. The target's pain from a unit of sanction, net of the
compliance cost of buying it off, is then at least its valuation of the political
state. The nine
equalities suffice for an exact threshold, and \Cref{prop:exact-thresholds}
gives exact thresholds without them, as at the partial-adjustment example.

\central{The sender's patience strictly raises the bounds on its own marginal
return to sanctioning, but with these nine equalities it does not move the
target's threshold.} Sender patience changes the level of the sender's response, and so
the compliance that relief requires, while the target's comparison turns on its
cost, pain, state valuation and discount factor.
On the class as a whole, sender patience also moves the curvature and Lipschitz
bounds in both sufficient conditions. The same structure brackets the sender. If the
upper substitution rate $k_{\mathrm{hi}}$ is below the lower bound
$u_{\mathrm{lo}}+\beta_Sg_Sd_{\mathrm{lo}}$ on the sender's dynamic marginal
return to sanctioning, and a larger concession is never worth more to the sender
at zero sanction, a positive sanction remains in every equilibrium whose
compliance does not rise with the state, full compliance included. If the lower
substitution rate $k_{\mathrm{lo}}$ is at least the upper bound $A$ on that return
(the lifting condition), full compliance ends the sanction.
\citet[Lemma~2, Proposition~1]{AcemogluJensen2015} give robust discount-factor
comparative statics for a single agent's dynamic program and for aggregates of
price-taking agents. Ours concern one government's patience in a two-player
game. A myopic against a forward-looking target is thus one cut of a continuous
comparison. Along it the constructed equilibrium switches from relief to
non-compliance as the target's discount factor rises from the sufficient condition
for relief to that for non-compliance. 

\central{Patience, like every other target parameter, can change whether an
economy has a full-non-compliance equilibrium, and leaves the path of such an
equilibrium unchanged.} Take two
economies that share the transition, the sender's discount factor and the
sender's stage payoff at full non-compliance ($\Us(x,0,s)$ on the square). In any full-non-compliance equilibria they have the same sender value
and the same state and sanction paths from every start, whatever the target's
payoff and discount factor and the sender's payoff away from non-compliance
(\Cref{prop:branch-invariance}).

\subsection{Welfare: the interest-group polity in full}
\label{sec:moved-main-2026-10-13:welfare}

The interest-group polity's accounting identity, the planner results and the
rescaling argument.

\central{In the interest-group polity of \Cref{app:fourgroup}, political weights
account exactly for the gap between the sum of the two governments' objectives and
aggregate welfare.} In the tradition of \citet{Kaempfer1988}, each country has two
interest groups. The sender has a pro-sanction group and a commercial lobby harmed
by sanctions. The target has a regime-aligned group and an opposition. Let
$W_g(x,a,s)$ be group $g$'s welfare and $\theta_g$ the political weight its
government places on it \citep{GrossmanHelpman1994}. \citet{Becker1983} supplies
the pressure contest behind the weights. Aggregate welfare $W$
sets every political weight to one and keeps the resource terms that enter the
government objectives at unit weight. Four sign conditions on the coefficients
deliver the game's stage assumptions. The sender's sanctioning cost is strictly
convex, a higher political state raises the sender's marginal value of
sanctioning, and compliance does not raise it. On the target side, the
regime-aligned group's gain in political
effectiveness under sanctions, weighted by its valuation of the political state,
outweighs the opposition's gain, weighted by the opposition's loss from the state.
For every coefficient vector, whether or not these conditions hold,
\begin{equation}
\Us(x,a,s)+\Ut(x,a,s)
=W(x,a,s)
+\sum_{g\in G}
(\theta_g-1)W_g(x,a,s),
\label{eq:planner-distortion}
\end{equation}
where $G$ is the set of the four groups. The identity is accounting at common
actions.

\central{At one example economy of this polity (\Cref{sec:model:examples}),
equilibrium play follows the corner path of full sanctions against full non-compliance
from every start, while the zero-sanction path attains the planner's optimum over
all feasible paths.} The planner's value exceeds equilibrium welfare at every
start and every planner discount factor (\Cref{prop:planner-optimum-lower-bound}).

\central{The zero-sanction path attains the planner's optimum in every
interest-group polity and under every transition satisfying \textup{(F2)--(F3)},
provided a sufficient condition holds: the first unit of sanction must not pay
socially.} Write $g_J$
and $q_K$ for the pro-sanction group's gain and the commercial lobby's loss per
unit of sanction. The target's direct cost per unit is $d_s$, and $B$ is the
regime-aligned group's valuation of the political state less the opposition's.
The zero-sanction path with the best constant compliance attains the optimum
from every start if $B\le0$ and $g_J-q_K\le d_s$. It also does so
if $B\ge0$ and
\[
g_J-q_K+\frac{\beta BL}{1-\beta\lambda}\le d_s ,
\]
where $L$ bounds the rise in next period's state per unit of sanction and
$\lambda$, with $\beta\lambda<1$, bounds the slope of the zero-sanction section
$\F(\cdot,0)$ (\Cref{prop:planner-optimum-lower-bound}). The condition says that
the sender-side groups' net gain plus the planner's discounted value of the higher
state the sanction leaves behind does not cover the target's direct cost.

\central{The comparison of the two paths extends beyond the example economy as
well.} For every vector of group coefficients and every transition satisfying
\textup{(F2)--(F3)}, the welfare difference between the zero-sanction path at full
non-compliance and the corner path is the capitalized long-run gap plus $B$ times a
discounted residual. The residual is one minus the distance between the two paths, which starts at
one. The sender's sanctioning cost is $k_S$, so the resource cost of a full
sanction to both countries is $k_S+d_s$. Under \textup{(F1)--(F5)} the difference
is positive at every planner discount if and only if this cost exceeds the
sender-side groups' net benefit $g_J-q_K$ and is at least that benefit plus $B$. The closed form
under the partial-adjustment law of \Cref{sec:model:examples} and the stationary
comparison are in \Cref{app:fourgroup}
(\Cref{prop:discounted-stationary-welfare,prop:affine-welfare-gap}). For other
transitions, the appendix bounds the gap below in the transition's own rates, more
sharply given a lower bound on its state slope along the zero-sanction and
full-sanction sections.

\central{The example shows that politically sustained sanctions \emph{can} be
socially wasteful, and neither the four sign conditions nor the game determines
the sign of the welfare gap.}
Multiply the pro-sanction group's per-unit benefit $g_J$ by $\phi>0$ and divide its
political-weight schedule by $\phi$. Both government payoffs are unchanged. So
under every transition the rescaled polity still satisfies the four sign
conditions and has the same game and equilibria. The planner, who counts that
group's benefit at unit weight, sees the gap between the zero-sanction path and
the corner path fall by $(\phi-1)g_J/(1-\beta)$. Hence, for every polity whose
pro-sanction group gains from sanctions, some rescaling makes the corner path
strictly better for the planner than the zero-sanction path, at every planner
discount, every start and every transition satisfying \textup{(F2)--(F3)}. Where
that group loses from sanctions, some rescaling makes the corner path strictly
worse (\Cref{sec:fourgroup:planner}). Each statement goes one way. Where the group
gains, the game never shows the zero-sanction path to be at least as good for the
planner; where it loses, never the reverse.

\central{Two polities next to the example economy reverse the welfare ranking,
the first with the same game and equilibrium.} The first rescales the example
economy in this way and keeps its equilibrium, hence its play of the corner path.
The second differs from it only in a lower direct sanction cost. The first fails
the sufficient condition above at every discount, and the second above a
threshold discount. For the first at every discount, and for the second above a
higher threshold, the corner path is worth strictly more to the planner than the
zero-sanction path. There the planner's optimum is not the zero-sanction path
(\Cref{sec:fourgroup:planner}).

\central{Maintained restrictions on the microfoundation narrow this
non-identification: they partially identify the ranking of cases $10$ and $11$, in
which the target maintains non-compliance.} With the target side of the microfoundation and the
commercial lobby's exposure held fixed, and the pro-sanction group's political
weight confined to a band, the welfare gap between the two constructed equilibria
ranges over an attained closed interval. A uniform strict ranking is identified
exactly when that interval excludes zero. Both strict rankings occur when its
endpoints have opposite signs. The exposure and the target-side inputs are
maintained from outside the governments' game. A weight
band alone need not identify the sign (\Cref{prop:welfare-interval}).
